Theorem \ref{thm:measure-for-universal} has potential to prove (or disprove) the universality of $\CGrid$ as a resource. However, it will be challenging to find a function $E$ that fails (satisfies Theorem \ref{thm:measure-for-universal}) for $\CGrid$ yet succeeds for $\Grid$. We formalize this attempt in the next subsection.

\subsection{Local Unitary vs. Local Clifford}

\begin{definition}
If $\ket{G_1} \geq_\LOCC \ket{G_2}$ and $\ket{G_2} \geq_\LOCC \ket{G_1}$ then $E(\ket{G_1}) = E(\ket{G_2})$ for any entanglement monotone $E$. An $E$ is called \term{entanglement monotone} if it does not increase, on average, for any $\LO$.
\end{definition}

\begin{definition}
$n$-qubit states $\ket\psi$ and $\ket\phi$ are LU-equivalent if there is a local unitary $U = \otimes_{i=1}^n U_i$ such that $\ket\psi = U \ket\phi$. We say $\ket\psi =_\LU \ket\phi$.
\end{definition}

\begin{definition}
The \term{one-qubit Clifford group} are the 1-qubit unitary matrices generated by $\langle H, S \rangle$. They are the stabilizers of the \term{Pauli group} $\cP = \langle \{-1,1,-i,i\} \rangle \times \{I,X,Y,Z\} \rangle$.
\end{definition}

\begin{definition}
$n$-qubit states $\ket\psi$ and $\ket\phi$ are LC-equivalent if there is a local  $U = \otimes_{i=1}^n U_i$ such that all $U_i$ belong to the one-qubit Clifford group and $\ket\psi = U \ket\phi$. We say $\ket\psi =_\LC \ket\phi$.
\end{definition}

\begin{theorem}
The LC-equivalence of graph states $\ket{G_1}$ and $\ket{G_2}$ is exactly determined by local complementation [cite:need]. 
\end{theorem}

\begin{definition}
Let $G = (V,E)$ be a graph. The local complementation around a vertex $v \in V$ is done by replacing the induced subgraph (the vertices adjacent to $v$, and any edges between these vertices) of $v$'s neighborhood with its graph complement.
\end{definition}

\begin{theorem}
$\ket{G_1} =_\LC \ket{G_2}$ if and only if there is a sequence of local complementations from $G_1$ to $G_2$. Furthermore, there exists an efficient algorithm to recognize whether graphs are local clifford equivalent \cite{lc-equiv-alg}.
\end{theorem}